\documentclass[12pt,a4paper]{article}
\usepackage[utf8]{inputenc}
\usepackage[T2A]{fontenc}
\usepackage[russian]{babel}
\usepackage{indentfirst}
\usepackage{graphicx}
\usepackage{geometry}
\usepackage{titlesec}
\usepackage{hyperref}
\usepackage{longtable}
\usepackage{array}
\usepackage{float}
\usepackage{xcolor}
\usepackage{listings}
\usepackage{tikz}
\usetikzlibrary{arrows.meta,positioning,fit}

\geometry{left=2cm,right=1.5cm,top=2cm,bottom=2cm}
\setlength{\parindent}{1.25cm}
\setlength{\parskip}{0.2em}

\hypersetup{
    colorlinks=true,
    linkcolor=black,
    urlcolor=blue
}

\lstdefinestyle{sqlstyle}{
    language=SQL,
    basicstyle=\small\ttfamily,
    keywordstyle=\color{blue}\bfseries,
    commentstyle=\color{gray},
    stringstyle=\color{purple},
    breaklines=true,
    frame=single,
    showstringspaces=false,
    columns=fullflexible
}

\lstset{style=sqlstyle}

\titleformat{\section}{\normalfont\Large\bfseries}{\thesection}{1em}{}
\titleformat{\subsection}{\normalfont\large\bfseries}{\thesubsection}{1em}{}

\begin{document}

\begin{titlepage}
    \centering
    \footnotesize
    Министерство цифрового развития, связи и массовых коммуникаций РФ\\
    Федеральное государственное образовательное учреждение высшего образования\\
    «Сибирский государственный университет телекоммуникаций и информатики»\\
    (СибГУТИ)

    \vspace{4cm}

    \Large \textbf{ОТЧЕТ}\\
    \large по расчетно-графической работе\\
    по дисциплине «Базы данных»\\
    \textbf{«Проектирование и реализация базы данных»}

    \vspace{1cm}

    \large
    \textbf{Тема:}\\
    «Система учета сетевой топологии и передачи трафика в локальной сети»

    \vspace{4cm}

    \hfill
    \begin{minipage}{0.48\textwidth}
        \textbf{Выполнили:}\\
        студенты группы ИКС-433\\
        Криволапов Н.\\
        Добромилов А.

        \vspace{1cm}

        \textbf{Проверил:}\\
        Брагин К.И.
    \end{minipage}

    \vfill

    Новосибирск 2026
\end{titlepage}

\tableofcontents
\newpage

\section{Цель работы}

Целью расчетно-графической работы является проектирование и реализация реляционной базы данных в СУБД PostgreSQL для предметной области, связанной с телекоммуникационными технологиями и IT-инфраструктурой.

В рамках работы требуется:
\begin{itemize}
    \item описать выбранную предметную область;
    \item спроектировать структуру базы данных;
    \item реализовать таблицы, первичные и внешние ключи;
    \item добавить ограничения целостности, индексы, представления, функции и триггеры;
    \item подготовить тестовое наполнение базы данных;
    \item разработать набор функциональных и аналитических SQL-запросов;
    \item обеспечить запуск PostgreSQL через Docker Compose.
\end{itemize}

\section{Описание предметной области}

В качестве предметной области выбрана система учета сетевой топологии и передачи трафика в локальной сети. Такая система может использоваться как часть NMS (Network Management System) для хранения информации об узлах сети, коммутаторах, портах, подключениях, MAC-таблице и событиях передачи Ethernet-кадров.

В локальной сети рабочие станции и серверы подключаются к портам коммутаторов. Коммутатор выполняет пересылку Ethernet-кадров между портами, используя MAC-таблицу. MAC-таблица содержит сведения о том, какой MAC-адрес был замечен на каком порту. При передаче кадра система фиксирует отправителя, получателя, MAC-адреса, тип кадра, размер полезной нагрузки и путь прохождения кадра через сетевое оборудование.

Основные автоматизируемые задачи:
\begin{itemize}
    \item ведение реестра сетевых хостов;
    \item учет коммутаторов и их портов;
    \item учет активных подключений хостов к портам;
    \item хранение MAC-таблицы коммутаторов;
    \item регистрация фактов передачи кадров;
    \item хранение пути прохождения кадров через порты коммутаторов;
    \item построение аналитических отчетов по топологии и трафику.
\end{itemize}

\section{Логическая модель базы данных}

База данных содержит 7 основных таблиц:
\begin{enumerate}
    \item \texttt{hosts} --- хосты локальной сети;
    \item \texttt{switches} --- сетевые коммутаторы;
    \item \texttt{switch\_ports} --- порты коммутаторов;
    \item \texttt{host\_connections} --- подключения хостов к портам;
    \item \texttt{mac\_table} --- MAC-таблица коммутаторов;
    \item \texttt{frames} --- Ethernet-кадры;
    \item \texttt{frame\_path} --- путь прохождения кадров.
\end{enumerate}

\subsection{ER-диаграмма}

\begin{figure}[H]
\centering
\resizebox{\textwidth}{!}{
\begin{tikzpicture}[
    entity/.style={draw, rounded corners, align=left, inner sep=6pt, font=\small, fill=gray!5},
    rel/.style={-{Latex[length=3mm]}, thick}
]
    \node[entity] (hosts) {
        \textbf{hosts}\\
        \underline{host\_id}\\
        hostname\\
        ip\_address\\
        mac\_address\\
        status\\
        created\_at
    };

    \node[entity, right=3.2cm of hosts] (connections) {
        \textbf{host\_connections}\\
        \underline{connection\_id}\\
        host\_id\\
        port\_id\\
        connected\_at\\
        disconnected\_at
    };

    \node[entity, right=3.2cm of connections] (ports) {
        \textbf{switch\_ports}\\
        \underline{port\_id}\\
        switch\_id\\
        port\_number\\
        status\\
        speed\_mbps
    };

    \node[entity, right=3.2cm of ports] (switches) {
        \textbf{switches}\\
        \underline{switch\_id}\\
        name\\
        model\\
        ports\_count\\
        location\\
        created\_at
    };

    \node[entity, below=2.3cm of ports] (mac) {
        \textbf{mac\_table}\\
        \underline{entry\_id}\\
        port\_id\\
        mac\_address\\
        learned\_at\\
        last\_seen\_at\\
        is\_static
    };

    \node[entity, below=2.7cm of hosts] (frames) {
        \textbf{frames}\\
        \underline{frame\_id}\\
        src\_host\_id\\
        dst\_host\_id\\
        src\_mac\\
        dst\_mac\\
        ether\_type\\
        payload\_size\\
        sent\_at\\
        delivery\_status
    };

    \node[entity, below=2.7cm of connections] (path) {
        \textbf{frame\_path}\\
        \underline{path\_id}\\
        frame\_id\\
        switch\_id\\
        in\_port\_id\\
        out\_port\_id\\
        step\_number\\
        processed\_at
    };

    \draw[rel] (connections) -- node[above, font=\scriptsize]{N:1 host\_id} (hosts);
    \draw[rel] (connections) -- node[above, font=\scriptsize]{N:1 port\_id} (ports);
    \draw[rel] (ports) -- node[above, font=\scriptsize]{N:1 switch\_id} (switches);
    \draw[rel] (mac) -- node[right, font=\scriptsize]{N:1 port\_id} (ports);
    \draw[rel] (frames) -- node[left, font=\scriptsize]{N:1 src/dst host} (hosts);
    \draw[rel] (path) -- node[above, font=\scriptsize]{N:1 frame\_id} (frames);
    \draw[rel] (path) -- node[right, font=\scriptsize]{N:1 switch\_id} (switches);
    \draw[rel] (path) -- node[right, font=\scriptsize]{N:1 in/out ports} (ports);
\end{tikzpicture}
}
\caption{ER-диаграмма базы данных с кратностями связей}
\end{figure}

\subsection{Связи между таблицами}

\begin{longtable}{|p{4.7cm}|p{4.7cm}|p{1.6cm}|p{4.1cm}|}
\hline
\textbf{Дочерняя таблица и внешний ключ} & \textbf{Родительская таблица и первичный ключ} & \textbf{Кратность} & \textbf{Смысл связи}\\
\hline
\texttt{host\_connections.host\_id} & \texttt{hosts.host\_id} & N:1 & Один хост может иметь несколько записей истории подключений.\\
\hline
\texttt{host\_connections.port\_id} & \texttt{switch\_ports.port\_id} & N:1 & Один порт может встречаться в истории подключений.\\
\hline
\texttt{switch\_ports.switch\_id} & \texttt{switches.switch\_id} & N:1 & Один коммутатор содержит много портов.\\
\hline
\texttt{mac\_table.port\_id} & \texttt{switch\_ports.port\_id} & N:1 & На одном порту может быть изучено несколько MAC-адресов.\\
\hline
\texttt{frames.src\_host\_id} & \texttt{hosts.host\_id} & N:1 & Хост может быть отправителем многих кадров.\\
\hline
\texttt{frames.dst\_host\_id} & \texttt{hosts.host\_id} & N:1 & Хост может быть получателем многих кадров.\\
\hline
\texttt{frame\_path.frame\_id} & \texttt{frames.frame\_id} & N:1 & Один кадр может иметь несколько шагов пути.\\
\hline
\texttt{frame\_path.switch\_id} & \texttt{switches.switch\_id} & N:1 & Один коммутатор может участвовать во многих шагах прохождения кадров.\\
\hline
\texttt{frame\_path.in\_port\_id} & \texttt{switch\_ports.port\_id} & N:1 & Порт может быть входным для многих шагов прохождения кадров.\\
\hline
\texttt{frame\_path.out\_port\_id} & \texttt{switch\_ports.port\_id} & N:1 & Порт может быть выходным для многих шагов прохождения кадров.\\
\hline
\end{longtable}

\subsection{Описание таблиц}

\begin{longtable}{|p{3.2cm}|p{4.2cm}|p{7.5cm}|}
\hline
\textbf{Таблица} & \textbf{Назначение} & \textbf{Ключевые поля}\\
\hline
\texttt{hosts} & Хранит конечные устройства сети & \texttt{host\_id}, \texttt{hostname}, \texttt{ip\_address}, \texttt{mac\_address}, \texttt{status}\\
\hline
\texttt{switches} & Хранит сведения о коммутаторах & \texttt{switch\_id}, \texttt{name}, \texttt{model}, \texttt{ports\_count}, \texttt{location}\\
\hline
\texttt{switch\_ports} & Описывает физические порты коммутаторов & \texttt{port\_id}, \texttt{switch\_id}, \texttt{port\_number}, \texttt{status}, \texttt{speed\_mbps}\\
\hline
\texttt{host\_connections} & Фиксирует подключение хоста к порту & \texttt{connection\_id}, \texttt{host\_id}, \texttt{port\_id}, \texttt{connected\_at}, \texttt{disconnected\_at}\\
\hline
\texttt{mac\_table} & Хранит MAC-таблицу коммутаторов & \texttt{entry\_id}, \texttt{port\_id}, \texttt{mac\_address}, \texttt{learned\_at}, \texttt{last\_seen\_at}\\
\hline
\texttt{frames} & Хранит факты передачи Ethernet-кадров & \texttt{frame\_id}, \texttt{src\_host\_id}, \texttt{dst\_host\_id}, \texttt{src\_mac}, \texttt{dst\_mac}, \texttt{payload\_size}\\
\hline
\texttt{frame\_path} & Хранит путь прохождения кадра через коммутаторы & \texttt{path\_id}, \texttt{frame\_id}, \texttt{switch\_id}, \texttt{in\_port\_id}, \texttt{out\_port\_id}, \texttt{step\_number}\\
\hline
\end{longtable}

\section{Реализация структуры базы данных}

Структура таблиц реализована в файле \texttt{sql/001\_init.sql}. Для каждой сущности задан первичный ключ. Межтабличные связи реализованы с помощью внешних ключей.

\begin{lstlisting}[caption={Пример создания таблиц hosts и switches}]
CREATE TABLE hosts (
    host_id SERIAL PRIMARY KEY,
    hostname VARCHAR(100) NOT NULL,
    ip_address INET UNIQUE NOT NULL,
    mac_address MACADDR UNIQUE NOT NULL,
    status VARCHAR(20) NOT NULL,
    created_at TIMESTAMP NOT NULL DEFAULT CURRENT_TIMESTAMP
);

CREATE TABLE switches (
    switch_id SERIAL PRIMARY KEY,
    name VARCHAR(100) NOT NULL,
    model VARCHAR(100),
    ports_count INT NOT NULL,
    location VARCHAR(100),
    created_at TIMESTAMP NOT NULL DEFAULT CURRENT_TIMESTAMP
);
\end{lstlisting}

Пример связи «коммутатор --- порт»:

\begin{lstlisting}[caption={Порты коммутаторов}]
CREATE TABLE switch_ports (
    port_id SERIAL PRIMARY KEY,
    switch_id INT NOT NULL REFERENCES switches(switch_id),
    port_number INT NOT NULL,
    status VARCHAR(20) NOT NULL,
    speed_mbps INT NOT NULL,
    UNIQUE (switch_id, port_number)
);
\end{lstlisting}

\section{Ограничения целостности}

Для обеспечения корректности данных реализованы ограничения \texttt{NOT NULL}, \texttt{UNIQUE}, \texttt{PRIMARY KEY}, \texttt{FOREIGN KEY} и дополнительные \texttt{CHECK}-ограничения.

\begin{longtable}{|p{5cm}|p{9.5cm}|}
\hline
\textbf{Ограничение} & \textbf{Назначение}\\
\hline
\texttt{chk\_hosts\_status} & Разрешает только допустимые статусы хоста: \texttt{online}, \texttt{offline}, \texttt{maintenance}.\\
\hline
\texttt{chk\_switches\_ports\_count} & Запрещает создавать коммутатор с неположительным числом портов.\\
\hline
\texttt{chk\_switch\_ports\_status} & Ограничивает статус порта значениями \texttt{up}, \texttt{down}, \texttt{disabled}.\\
\hline
\texttt{chk\_switch\_ports\_speed} & Проверяет, что скорость порта больше нуля.\\
\hline
\texttt{chk\_frames\_hosts\_are\_different} & Запрещает отправку кадра от хоста самому себе.\\
\hline
\texttt{chk\_frames\_payload\_size} & Проверяет положительный размер полезной нагрузки кадра.\\
\hline
\texttt{chk\_frames\_delivery\_status} & Ограничивает статус доставки значениями \texttt{delivered}, \texttt{dropped}, \texttt{flooded}.\\
\hline
\end{longtable}

\section{Индексы}

Индексы используются для ускорения частых операций поиска, фильтрации, соединения таблиц и аналитики по времени.

\begin{longtable}{|p{6cm}|p{8.5cm}|}
\hline
\textbf{Индекс} & \textbf{Обоснование}\\
\hline
\texttt{idx\_hosts\_status} & Ускоряет выборку хостов по статусу.\\
\hline
\texttt{idx\_switch\_ports\_switch\_id} & Ускоряет получение портов конкретного коммутатора.\\
\hline
\texttt{idx\_host\_connections\_host\_id} & Ускоряет поиск подключений конкретного хоста.\\
\hline
\texttt{idx\_host\_connections\_port\_id} & Ускоряет поиск подключений по порту.\\
\hline
\texttt{idx\_mac\_table\_mac\_address\_unique} & Обеспечивает уникальность MAC-адреса в MAC-таблице и ускоряет поиск.\\
\hline
\texttt{idx\_frames\_src\_host\_id} & Ускоряет отчеты по отправителям кадров.\\
\hline
\texttt{idx\_frames\_dst\_host\_id} & Ускоряет отчеты по получателям кадров.\\
\hline
\texttt{idx\_frames\_sent\_at} & Ускоряет аналитику по временным интервалам.\\
\hline
\texttt{idx\_frame\_path\_frame\_id} & Ускоряет трассировку пути конкретного кадра.\\
\hline
\end{longtable}

Пример создания индексов:

\begin{lstlisting}[caption={Индексы для таблиц frames и frame\_path}]
CREATE INDEX IF NOT EXISTS idx_frames_src_host_id ON frames(src_host_id);
CREATE INDEX IF NOT EXISTS idx_frames_dst_host_id ON frames(dst_host_id);
CREATE INDEX IF NOT EXISTS idx_frames_sent_at ON frames(sent_at);
CREATE INDEX IF NOT EXISTS idx_frame_path_frame_id ON frame_path(frame_id);
\end{lstlisting}

\section{Представления}

В базе данных созданы представления, упрощающие выполнение типовых аналитических запросов.

\begin{longtable}{|p{5cm}|p{9.5cm}|}
\hline
\textbf{Представление} & \textbf{Назначение}\\
\hline
\texttt{v\_active\_connections} & Показывает активные подключения хостов к портам коммутаторов.\\
\hline
\texttt{v\_switch\_mac\_table} & Показывает MAC-таблицу с названиями коммутаторов и номерами портов.\\
\hline
\texttt{v\_frame\_trace} & Показывает путь прохождения кадров через коммутаторы.\\
\hline
\texttt{v\_host\_traffic\_stats} & Сводит статистику отправленных и полученных кадров по хостам.\\
\hline
\end{longtable}

\section{Функции и триггеры}

Для реализации серверной логики используются функции PostgreSQL на SQL и PL/pgSQL.

\begin{longtable}{|p{5.5cm}|p{9cm}|}
\hline
\textbf{Функция} & \textbf{Назначение}\\
\hline
\texttt{learn\_mac\_address(...)} & Добавляет или обновляет запись в MAC-таблице.\\
\hline
\texttt{register\_frame(...)} & Регистрирует кадр и автоматически создает запись о его пути.\\
\hline
\texttt{get\_frame\_trace(...)} & Возвращает трассировку конкретного кадра.\\
\hline
\texttt{get\_host\_traffic\_stats()} & Возвращает статистику трафика по хостам.\\
\hline
\texttt{get\_switch\_mac\_table(...)} & Возвращает MAC-таблицу указанного коммутатора.\\
\hline
\end{longtable}

Пример функции регистрации кадра:

\begin{lstlisting}[caption={Вызов функции регистрации кадра}]
SELECT register_frame(1, 2, 'IPv4', 128, 'delivered') AS new_frame_id;
\end{lstlisting}

В функции \texttt{register\_frame} база данных сама находит активные подключения отправителя и получателя, получает их MAC-адреса и порты, создает запись в \texttt{frames}, затем создает запись в \texttt{frame\_path}.

Триггеры:

\begin{longtable}{|p{5.5cm}|p{9cm}|}
\hline
\textbf{Триггер} & \textbf{Назначение}\\
\hline
\texttt{trg\_validate\_frame\_hosts} & Проверяет, что MAC-адреса кадра соответствуют хостам-отправителю и получателю.\\
\hline
\texttt{trg\_learn\_src\_mac\_from\_frame} & После вставки кадра обновляет MAC-таблицу по MAC-адресу отправителя.\\
\hline
\texttt{trg\_validate\_frame\_path\_ports} & Проверяет, что входной и выходной порты пути принадлежат указанному коммутатору.\\
\hline
\end{longtable}

\section{Тестовые данные}

Тестовые данные находятся в файле \texttt{sql/002\_seed.sql}. В базу добавлены:
\begin{itemize}
    \item три хоста: две рабочие станции и сервер;
    \item один коммутатор Core-Switch-01;
    \item три порта коммутатора;
    \item активные подключения двух рабочих станций;
    \item начальная MAC-таблица;
    \item пример переданного кадра и его путь.
\end{itemize}

Пример наполнения:

\begin{lstlisting}[caption={Добавление хостов}]
INSERT INTO hosts (hostname, ip_address, mac_address, status) VALUES
('Workstation-01', '192.168.1.10', '00:50:56:c0:00:01', 'online'),
('Workstation-02', '192.168.1.11', '00:50:56:c0:00:02', 'online'),
('Server-Main', '192.168.1.200', '00:50:56:c0:00:ff', 'online');
\end{lstlisting}

\section{Функциональные SQL-запросы}

Файл \texttt{sql/queries.sql} содержит 20 демонстрационных запросов. Они покрывают простые выборки, фильтрацию, соединения таблиц, агрегацию, \texttt{HAVING}, CTE, сортировку, ограничение результата, временную аналитику и оконные функции. Ниже приведены основные типы запросов.

\subsection{Активные подключения}

\begin{lstlisting}
SELECT *
FROM v_active_connections
ORDER BY switch_id, port_number;
\end{lstlisting}

Запрос показывает, какие хосты подключены к каким портам коммутатора.

\subsection{MAC-таблица коммутатора}

\begin{lstlisting}
SELECT *
FROM get_switch_mac_table('Core-Switch-01');
\end{lstlisting}

Запрос возвращает MAC-таблицу указанного коммутатора.

\subsection{Трассировка кадра}

\begin{lstlisting}
SELECT *
FROM get_frame_trace(1);
\end{lstlisting}

Запрос показывает, через какой коммутатор и какие порты прошел конкретный кадр.

\subsection{Статистика трафика по хостам}

\begin{lstlisting}
SELECT *
FROM get_host_traffic_stats();
\end{lstlisting}

Запрос показывает количество и объем отправленных и полученных кадров по каждому хосту.

\subsection{Запрос с HAVING}

\begin{lstlisting}
SELECT
    h.hostname,
    COUNT(f.frame_id) AS frames_sent,
    SUM(f.payload_size) AS bytes_sent
FROM hosts h
JOIN frames f ON f.src_host_id = h.host_id
GROUP BY h.hostname
HAVING COUNT(f.frame_id) > 1
ORDER BY frames_sent DESC;
\end{lstlisting}

Запрос выбирает только тех хостов, которые отправили больше одного кадра.

\subsection{Запрос с CTE и LIMIT}

\begin{lstlisting}
WITH traffic_pairs AS (
    SELECT
        src.hostname AS src_hostname,
        dst.hostname AS dst_hostname,
        COUNT(*) AS frames_count,
        SUM(f.payload_size) AS total_bytes
    FROM frames f
    JOIN hosts src ON src.host_id = f.src_host_id
    JOIN hosts dst ON dst.host_id = f.dst_host_id
    GROUP BY src.hostname, dst.hostname
)
SELECT *
FROM traffic_pairs
ORDER BY frames_count DESC, total_bytes DESC
LIMIT 3;
\end{lstlisting}

Запрос строит рейтинг наиболее активных направлений передачи кадров.

\subsection{Дополнительные запросы 13--20}

Для расширения аналитического набора добавлены запросы, показывающие свободные ресурсы сети, распределение трафика, последние события и оконную нумерацию кадров.

\begin{longtable}{|p{1.3cm}|p{7cm}|p{6cm}|}
\hline
\textbf{Номер} & \textbf{Назначение запроса} & \textbf{Используемые конструкции}\\
\hline
13 & Хосты без активного подключения к коммутатору & \texttt{LEFT JOIN}, \texttt{WHERE IS NULL}\\
\hline
14 & Свободные порты коммутаторов & \texttt{JOIN}, \texttt{LEFT JOIN}, фильтрация\\
\hline
15 & Количество кадров по статусам доставки & \texttt{GROUP BY}, \texttt{COUNT}, \texttt{SUM}\\
\hline
16 & Динамика трафика по дням & \texttt{DATE}, \texttt{GROUP BY}, \texttt{AVG}\\
\hline
17 & Последние 5 переданных кадров с именами хостов & \texttt{JOIN}, \texttt{ORDER BY}, \texttt{LIMIT}\\
\hline
18 & Утилизация портов коммутаторов & \texttt{GROUP BY}, агрегаты, вычисляемые поля\\
\hline
19 & MAC-адреса на портах без активного подключения & \texttt{LEFT JOIN}, \texttt{WHERE IS NULL}\\
\hline
20 & Нумерация кадров по каждому отправителю & \texttt{ROW\_NUMBER()}, \texttt{PARTITION BY}\\
\hline
\end{longtable}

\section{Контейнеризация}

Для запуска PostgreSQL используется Docker Compose. Конфигурация находится в файле \texttt{docker-compose.yml}.

\begin{lstlisting}[caption={Фрагмент docker-compose.yml}]
services:
  postgres:
    image: postgres:16
    container_name: lan_traffic_postgres
    environment:
      POSTGRES_USER: postgres
      POSTGRES_PASSWORD: postgres
      POSTGRES_DB: lan_traffic_db
    ports:
      - "5432:5432"
    volumes:
      - postgres_data:/var/lib/postgresql/data
      - ./sql/001_init.sql:/docker-entrypoint-initdb.d/001_init.sql:ro
      - ./sql/002_seed.sql:/docker-entrypoint-initdb.d/002_seed.sql:ro
      - ./sql/003_db_objects.sql:/docker-entrypoint-initdb.d/003_db_objects.sql:ro
\end{lstlisting}

Запуск проекта:

\begin{lstlisting}[language=bash]
docker compose up -d
\end{lstlisting}

Подключение к базе данных:

\begin{lstlisting}[language=bash]
docker exec -it lan_traffic_postgres psql -U postgres -d lan_traffic_db
\end{lstlisting}

Полная очистка контейнера и тома данных:

\begin{lstlisting}[language=bash]
docker compose down -v
\end{lstlisting}

\section{Проверка работоспособности}

При проверке базы данных были получены следующие результаты:
\begin{itemize}
    \item создано 7 таблиц;
    \item создано 4 представления;
    \item создано 8 функций;
    \item создано 3 триггера;
    \item создано 26 индексов, включая системные индексы первичных и уникальных ключей;
    \item демонстрационный файл \texttt{queries.sql} выполняется без ошибок.
\end{itemize}

Пример проверки списка таблиц:

\begin{lstlisting}[language=bash]
docker exec lan_traffic_postgres psql -U postgres -d lan_traffic_db -c "\dt"
\end{lstlisting}

Пример проверки функций:

\begin{lstlisting}[language=bash]
docker exec lan_traffic_postgres psql -U postgres -d lan_traffic_db -c "\df"
\end{lstlisting}

\section{Соответствие требованиям задания}

\begin{longtable}{|p{8cm}|p{6.5cm}|}
\hline
\textbf{Требование} & \textbf{Выполнение}\\
\hline
Не менее 5 логически связанных таблиц & Реализовано 7 таблиц.\\
\hline
Первичные и внешние ключи & Реализованы во всех основных таблицах.\\
\hline
Тестовое наполнение & Реализовано в \texttt{sql/002\_seed.sql}.\\
\hline
Не менее 10 SQL-запросов & Реализовано 20 запросов в \texttt{sql/queries.sql}.\\
\hline
Ограничения \texttt{NOT NULL}, \texttt{UNIQUE}, \texttt{CHECK} & Реализованы в схеме и файле \texttt{003\_db\_objects.sql}.\\
\hline
Индексы & Реализованы индексы для связей, фильтрации и аналитики.\\
\hline
HAVING, CTE, ORDER BY LIMIT & Реализованы в демонстрационных запросах.\\
\hline
Docker Compose & Реализован запуск PostgreSQL через \texttt{docker-compose.yml}.\\
\hline
PL/pgSQL функции и триггеры & Реализованы функции и триггеры серверной логики.\\
\hline
\end{longtable}

\section{Заключение}

В ходе выполнения расчетно-графической работы была спроектирована и реализована база данных для учета сетевой топологии и передачи трафика в локальной сети. База данных позволяет хранить сведения о хостах, коммутаторах, портах, подключениях, MAC-таблице, Ethernet-кадрах и пути их прохождения.

Реализованная схема содержит ограничения целостности, индексы для ускорения запросов, представления для аналитики, функции на SQL и PL/pgSQL, а также триггеры для контроля бизнес-правил и автоматического обновления MAC-таблицы. Для удобного запуска и проверки используется Docker Compose с автоматической инициализацией PostgreSQL.

Разработанная база данных соответствует требованиям задания и может быть использована как учебная модель системы мониторинга и учета событий в локальной сети.

\end{document}
